\section{A normalized Mellin formula for the height-one germs}
\label{sec:laurent}

This section gives a self-contained normalization of the coefficient
data used later. The reciprocal-Gamma singular coefficients and
height-one series are established objects in the literature
\cite{Young2023,Young2024} and in the incoming Barnes report
\cite{Barnes}. We record the precise integral form so that the fifth
power and the ordinary Gamma integral share a proved convention.

For $v>0$, put
\begin{equation}\label{eq:Vdef}
 B(v)=-\log(1-e^{-v}),\qquad
 V_d(v)=B(v)^d-\boldsymbol1_{(0,1)}(v)(-\log v)^d.
\end{equation}
Near zero, $B(v)=-\log v+O(v)$, and consequently
\[
 V_d(v)=O\bigl(v(1+|\log v|^{d-1})\bigr).
\]
At infinity it decays exponentially. Thus
$\int_0^\infty v^{u-1}V_d(v)\,dv$ is holomorphic for
$\re u>-1$; each derivative on a compact subset is obtained by
inserting the corresponding power of $\log v$.

\begin{proposition}[Mellin normalization]\label{prop:mellin}
For $d\ge1$, initially for $\re u>0$ and then meromorphically for
$\re u>-1$,
\begin{equation}\label{eq:normalizedmellin}
 \boxed{\displaystyle
 \Gamma(1+u)Z_d(1+u,\{1\}^{d-1})
   =\frac1{u^d}+\frac{u}{d!}
       \int_0^\infty v^{u-1}V_d(v)\,dv.}
\end{equation}
In particular the complete principal part and the constant term are
\begin{equation}\label{eq:principal}
 Z_d(1+u,\{1\}^{d-1})
 =\sum_{j=0}^d c_j u^{j-d}+\lambda_du+O(u^2),
\end{equation}
where
\begin{equation}\label{eq:lambda_integral}
 \lambda_d=c_{d+1}+\frac1{d!}\int_0^\infty\frac{V_d(v)}v\,dv.
\end{equation}
\end{proposition}
\begin{proof}
The generating identity
\[
 \Li_{\{1\}^{d-1}}(y)=\frac{(-\log(1-y))^{d-1}}{(d-1)!}
\]
follows either from strict harmonic generating functions or by
repeated integration of $1/(1-y)$. Termwise Mellin integration gives
\[
 \Gamma(1+u)Z_d(1+u,\{1\}^{d-1})
  =\frac1{(d-1)!}\int_0^\infty
       \frac{v^u B(v)^{d-1}}{e^v-1}\,dv.
\]
For $\re u>0$ the integral is absolutely convergent. Since
$B'(v)=-1/(e^v-1)$, integration by parts, with both boundary terms
zero, turns it into
\[
 \frac{u}{d!}\int_0^\infty v^{u-1}B(v)^d\,dv.
\]
The subtracted comparison integral is exactly
$\int_0^1v^{u-1}(-\log v)^d\,dv=d!/u^{d+1}$.
The estimates preceding the proposition justify continuation and all
coefficient extractions. Multiplication by the Taylor series
\eqref{eq:cdef} proves the remaining formulas.
\end{proof}

\begin{corollary}[Every regular coefficient and spectral derivative]
\label{cor:allcoefficients}
Let
\[
 I_{d,r}=\int_0^\infty \frac{(\log v)^r V_d(v)}v\,dv
 \qquad(d\ge1,\ r\ge0).
\]
For every $j\ge1$,
\begin{equation}\label{eq:alllaurent}
 \lambda_{d,j}=c_{d+j}
  +\frac1{d!}\sum_{r=0}^{j-1}
       c_{j-1-r}\frac{I_{d,r}}{r!}.
\end{equation}
If $\Lambda_d(u)=\Gamma(1+u)Z_d(1+u,\{1\}^{d-1})-u^{-d}$,
then
\begin{equation}\label{eq:normalizedderivatives}
 \Lambda_d^{(j)}(0)=\frac{j}{d!}I_{d,j-1}\qquad(j\ge1).
\end{equation}
These are ordinary convergent logarithmic Mellin integrals. At depth
one, $\lambda_{1,j}=(-1)^j\gamma_j/j!$ recovers the usual Stieltjes
derivative convention.
\end{corollary}
\begin{proof}
Expand the holomorphic integral in \eqref{eq:normalizedmellin} into
its Taylor series and multiply by $\sum c_ju^j$. The estimates are
uniform on every closed disc $|u|\le\rho<1$, so the operations are
analytic coefficient extractions rather than formal regularizations.
\end{proof}

Formula~\eqref{eq:alllaurent} is useful for reproducibility: coefficients
of different depths can be evaluated in one fixed scheme, with no
matching of independently chosen subtraction constants. It is also a
direct route to higher spectral derivatives whenever a new digamma
identity produces them.
